% Source: https://artofproblemsolving.com/community/c3h1754235_list_of_math_problems_i_need_help_with
% Stable identity: aops-topic-1754235-q03
One admissible colouring is shown.

\begin{asy}
usepackage("amsmath");
size(5cm);
defaultpen(fontsize(10pt));
for(int i=0;i<7;++i) draw(box((i,0),(i+1,1)),linewidth(0.8));
fill(box((0,0),(1,1)),gray(0.6));fill(box((3,0),(4,1)),gray(0.6));fill(box((5,0),(6,1)),gray(0.6));
for(int i=0;i<7;++i) draw(box((i,0),(i+1,1)),linewidth(0.8));
label("\mbox{}",(min(currentpicture,user=true)+max(currentpicture,user=true))/2,black); // keeps the text after the figure black when LaTeX typesets the labels inline
\end{asy}

Remove one uncoloured square from each of the $2$ gaps that must separate the $3$ coloured squares. This gives a bijection between admissible colourings of the $7$-strip and arbitrary choices of $3$ squares in a $5$-strip (put the separating squares back to recover the colouring). Hence the number of ways is $\dbinom53=\boxed{10}$.

Directly: if the coloured squares are in positions $a<b<c$ with $b\ge a+2$ and $c\ge b+2$, then $a<b-1<c-2$ are any three distinct positions among $1,\dots,5$.
